% Einheit 1 – Wurzelziehen und Lösbarkeit
\clearpage
\setcounter{aufgabe}{9}
\hypertarget{e1}{}\einheitenkopf{Einheit 1 von 4 · Wurzelziehen und Lösbarkeit}
\zweigzeile{Hier lernst du, Gleichungen wie $x^2 = c$ durch Wurzelziehen zu lösen und zu entscheiden, ob sie zwei, eine oder keine Lösung haben · neu oder schon bekannt – je nach Buch (an der Oberschule Klasse 9 oder 10) · P10 · baut auf: Quadrieren, Quadratwurzeln, lineare Gleichungen, Scheitelpunktform}

\begin{aufgabe}{Ich erkenne, ob eine Gleichung quadratisch oder linear ist. Kreuze an. Löse nichts.}
\begin{teile}
\teil $x^2 = 50$ \hfill \kreuz{quadratisch}\kreuz{linear}
\teil $4x - 1 = 19$ \hfill \kreuz{quadratisch}\kreuz{linear}
\teil $(x - 1)^2 = 25$ \hfill \kreuz{quadratisch}\kreuz{linear}
\teil $3\cdot(x + 2) = 18$ \hfill \kreuz{quadratisch}\kreuz{linear}
\teil $x\cdot x = 14$ \hfill \kreuz{quadratisch}\kreuz{linear}
\end{teile}
\end{aufgabe}

\begin{aufgabe}{Ich erkenne an der rechten Seite von $x^2 = c$, ob es zwei, eine oder keine Lösung gibt. Sieh nur die rechte Seite an und kreuze an. Rechne nichts.}
\begin{teile}
\teil $x^2 = 30$ \hfill \kreuz{positiv}\kreuz{null}\kreuz{negativ}\quad$\Rightarrow$\quad\kreuz{zwei}\kreuz{eine}\kreuz{keine}
\teil $x^2 = -1$ \hfill \kreuz{positiv}\kreuz{null}\kreuz{negativ}\quad$\Rightarrow$\quad\kreuz{zwei}\kreuz{eine}\kreuz{keine}
\teil $x^2 = 0$ \hfill \kreuz{positiv}\kreuz{null}\kreuz{negativ}\quad$\Rightarrow$\quad\kreuz{zwei}\kreuz{eine}\kreuz{keine}
\teil $x^2 = 0{,}1$ \hfill \kreuz{positiv}\kreuz{null}\kreuz{negativ}\quad$\Rightarrow$\quad\kreuz{zwei}\kreuz{eine}\kreuz{keine}
\teil $x^2 = -100$ \hfill \kreuz{positiv}\kreuz{null}\kreuz{negativ}\quad$\Rightarrow$\quad\kreuz{zwei}\kreuz{eine}\kreuz{keine}
\end{teile}
\end{aufgabe}

\begin{aufgabe}{Ich kann $x^2 = c$ durch Wurzelziehen lösen. Gib beide Lösungen an.}
\beispiel{x^2 &= 81 &&\mid \sqrt{\phantom{x}} \\ x_1 &= 9, \quad x_2 = -9}
\begin{gleichungsraster}[1]
\gl{x^2 = 25} & \gl{x^2 = 100} \\ \noalign{\vspace{6pt}}
\gl{x^2 = 144} & \gl{x^2 = 1} \\ \noalign{\vspace{6pt}}
\end{gleichungsraster}
\medskip
\begin{teile}
\teil Ein quadratischer Pausenhof hat den Flächeninhalt $225\,\text{m}^2$. Wie lang ist eine Seite? Welche Lösung der Gleichung passt nicht zur Frage?
\end{teile}
\medskip
Die Wurzel geht nicht auf: Lass sie stehen und gib zusätzlich Näherungswerte auf zwei Nachkommastellen an.
\begin{gleichungsraster}[2]
\gl{x^2 = 14} & \\ \noalign{\vspace{6pt}}
\end{gleichungsraster}
\medskip
\begin{teile}
\teil Begründe, warum $x^2 = -64$ keine Lösung hat.
\end{teile}
\end{aufgabe}

\begin{aufgabe}{Ich kann $x^2$ erst allein stellen und dann die Wurzel ziehen. Löse die Gleichung.}
\beispiel{x^2 - 7 &= 57 &&\mid +7 \\ x^2 &= 64 &&\mid \sqrt{\phantom{x}} \\ x_1 &= 8, \quad x_2 = -8}
\begin{gleichungsraster}[2]
\gl{x^2 + 5 = 30} & \gl{2x^2 = 98} \\ \noalign{\vspace{6pt}}
\gl[3]{2x^2 - 3 = 29} & \\ \noalign{\vspace{6pt}}
\end{gleichungsraster}
\medskip
\begin{teile}
\teil Eine zylinderförmige Dose soll $850\,\text{cm}^3$ fassen und $12\,\text{cm}$ hoch sein. Stelle $V = \pi\cdot r^2\cdot h$ nach $r^2$ um und berechne den Radius. Runde auf eine Nachkommastelle.
\end{teile}
\medskip
Stelle $x^2$ allein und entscheide dann: zwei, eine oder keine Lösung? Gib die Lösungen an.
\begin{gleichungsraster}[2]
\gl{4x^2 + 9 = 9} & \gl{2x^2 + 11 = 3} \\ \noalign{\vspace{6pt}}
\end{gleichungsraster}
\end{aufgabe}

\begin{aufgabe}{Ich kann eine Gleichung mit quadrierter Klammer lösen. Ziehe die Wurzel mit beiden Vorzeichen und rechne dann rückwärts.}
\beispiel{(x - 3)^2 &= 25 &&\mid \sqrt{\phantom{x}} \\ x - 3 &= 5 \quad\text{oder}\quad x - 3 = -5 &&\mid +3 \\ x_1 &= 8, \quad x_2 = -2}
\begin{gleichungsraster}[2]
\gl{(x - 1)^2 = 9} & \gl{(x - 5)^2 = 4} \\ \noalign{\vspace{6pt}}
\gl{(x + 2)^2 = 49} & \gl{(x + 6)^2 = 0} \\ \noalign{\vspace{6pt}}
\end{gleichungsraster}
\medskip
\begin{teile}
\teil Prüfe durch Einsetzen, ob $x = -5$ eine Lösung von $(x + 2)^2 = 9$ ist. \hfill \kreuz{(wA)}\kreuz{(fA)}
\end{teile}
\end{aufgabe}

\begin{aufgabe}{Ich kann am Graphen ablesen, wie viele Lösungen eine Gleichung hat. Die Parabel gehört zu $y = (x - 1)^2 - 3$. Die Lösungen von $(x - 1)^2 - 3 = c$ sind die $x$-Werte der Schnittpunkte der Parabel mit der waagerechten Geraden $y = c$. Lies ab, wie viele Lösungen es gibt, und gib sie an.}

\begin{ksys}[xmin=-3,xmax=5,ymin=-5,ymax=5,ablesen]
\parabel{1}{1}{-3}{}
\gerade[4.3]{0}{1}{y=1}
\gerade[2.8]{0}{-3}{y=-3}
\gerade[2.8]{0}{-4}{y=-4}
\end{ksys}

\begin{teile}
\teil $(x - 1)^2 - 3 = 1$ \hfill Anzahl: \leerfeld Lösungen: \leerfeld
\teil $(x - 1)^2 - 3 = -3$ \hfill Anzahl: \leerfeld Lösungen: \leerfeld
\teil $(x - 1)^2 - 3 = -4$ \hfill Anzahl: \leerfeld Lösungen: \leerfeld
\teil Gib eine Zahl $c$ an, für die $(x - 1)^2 - 3 = c$ keine Lösung hat. \hfill $c = $ \leerfeld
\end{teile}
\end{aufgabe}

\begin{aufgabe}{Ich kann Aussagen über die Lösungen quadratischer Gleichungen beurteilen.}
\begin{teile}
\teil Kreuze an, ob die Aussage wahr oder falsch ist.
\end{teile}
\sachtabelle{lcc}{Aussage & wahr & falsch}{$(x + 5)^2 = 0$ hat genau eine Lösung. & $\square$ & $\square$ \\ $3$ und $9$ sind Lösungen von $(x + 6)^2 = 9$. & $\square$ & $\square$}
\begin{teile}
\teil Gib eine quadratische Gleichung an, die keine Lösung hat. Begründe. (P10 2021)
\end{teile}
\end{aufgabe}

\begin{aufgabe}{Ich finde den Fehler beim Wurzelziehen. Lena löst $(x + 1)^2 = 81$ so:}
\rechnung{(x + 1)^2 &= 81 &&\mid \sqrt{\phantom{x}} \\ x + 1 &= 9 &&\mid -1 \\ x &= 8}
\begin{teile}
\teil Finde den Fehler, benenne ihn und korrigiere die Rechnung.
\end{teile}
\medskip
Löse selbst.
\begin{gleichungsraster}[2]
\gl{(x + 3)^2 = 64} & \\ \noalign{\vspace{6pt}}
\end{gleichungsraster}
\end{aufgabe}

\begin{aufgabe}{Ich kann begründen, warum $x^2 = c$ zwei, eine oder keine Lösung hat. Begründe ohne Rechnung.}
\begin{teile}
\teil Warum hat $x^2 = 40$ zwei Lösungen?
\teil Warum hat $x^2 = -40$ keine Lösung?
\teil Tim sagt: „$x^2 = 0$ hat keine Lösung, denn aus null kann man keine Wurzel ziehen.“ Hat Tim recht? Begründe.
\end{teile}
\end{aufgabe}
